
\documentclass{article}
\usepackage{colortbl}
\usepackage{makecell}
\usepackage{multirow}
\usepackage{supertabular}

\begin{document}

\newcounter{utterance}

\centering \large Interaction Transcript for game `cladder', experiment `full\_v1.5\_default', episode 8837 with refineonly{-}Qwen3.5{-}9B.
\vspace{24pt}

{ \footnotesize  \setcounter{utterance}{1}
\setlength{\tabcolsep}{0pt}
\begin{supertabular}{c@{$\;$}|p{.15\linewidth}@{}p{.15\linewidth}p{.15\linewidth}p{.15\linewidth}p{.15\linewidth}p{.15\linewidth}}
    \# & \multicolumn{2}{c}{Player} && \multicolumn{2}{c}{Game Master} \\
    \hline

    \theutterance \stepcounter{utterance}  
    & & & \multicolumn{4}{p{0.6\linewidth}}{
        \cellcolor[rgb]{0.9,0.9,0.9}{
            \makecell[{{p{\linewidth}}}]{
                \texttt{\tiny{[P1$\langle$GM]}}
                \texttt{You are an expert in causal inference. The following question is not a typical commonsense query, but rather a meticulously designed question created by a professor specializing in causal inference, intended to assess the students' mastery of the course content.} \\
\texttt{Imagine a self{-}contained, hypothetical world with only the following conditions, and without any unmentioned factors or causal relationships: Yupt has a direct effect on zupj and muvq. Xyfo has a direct effect on zupj and muvq. Zupj has a direct effect on muvq. Xyfo is unobserved. For those who are not yupt and are not zupj, the probability of muvq is 44\%. For those who are not yupt and are zupj, the probability of muvq is 51\%. For those who are yupt and are not zupj, the probability of muvq is 57\%. For those who are yupt and are zupj, the probability of muvq is 10\%. For those who are not yupt and are not xyfo, the probability of zupj is 42\%. For those who are not yupt and are xyfo, the probability of zupj is 87\%. For those who are yupt and are not xyfo, the probability of zupj is 26\%. For those who are yupt and are xyfo, the probability of zupj is 44\%. The overall probability of xyfo is 22\%. Does yupt negatively affect muvq through zupj?} \\
\texttt{Start your answer with 'Yes' or 'No', followed by additional reasoning or evidence to support your explanation.} \\
            }
        }
    }
     \\ \\

    \theutterance \stepcounter{utterance}  
    & \multicolumn{4}{p{0.6\linewidth}}{
        \cellcolor[rgb]{0.561,0.737,0.561}{
            \makecell[{{p{\linewidth}}}]{
                \texttt{\tiny{[P1$\rangle$GM]}}
                \texttt{No} \\
\\ 
\texttt{**Reasoning:**} \\
\\ 
\texttt{To determine if \$Y\_\{upt\}\$ negatively affects \$M\_\{uvq\}\$ through \$Z\_\{upj\}\$, we must calculate the **Natural Indirect Effect (NIE)**, which represents the change in the probability of \$M\_\{uvq\}\$ resulting from the change in \$Z\_\{upj\}\$ induced by changing \$Y\_\{upt\}\$ from 0 to 1, while holding \$Y\_\{upt\}\$ fixed at 0.} \\
\\ 
\texttt{The causal graph structure is:} \\
\texttt{\$X\_\{yfo\} \textbackslash{}to Z\_\{upj\} \textbackslash{}to M\_\{uvq\}\$} \\
\texttt{\$Y\_\{upt\} \textbackslash{}to Z\_\{upj\}\$} \\
\texttt{\$Y\_\{upt\} \textbackslash{}to M\_\{uvq\}\$} \\
\\ 
\texttt{**Step 1: Define the Counterfactual Probabilities**} \\
\texttt{We need to evaluate \$P(M\_\{uvq\} = 1 | Y\_\{upt\}=0, Z\_\{upj\} = z)\$ for the relevant values of \$z\$ (\$0\$ and \$1\$).} \\
\texttt{From the problem statement:} \\
\texttt{*   \$P(M=1 | Y=0, Z=0) = 44\textbackslash{}\%\$} \\
\texttt{*   \$P(M=1 | Y=0, Z=1) = 51\textbackslash{}\%\$} \\
\\ 
\texttt{**Step 2: Calculate the Indirect Effect (NIE)**} \\
\texttt{The NIE is calculated as the difference in the probability of \$M\_\{uvq\}\$ when the mediator \$Z\_\{upj\}\$ changes from its baseline state (\$Z=0\$) to its intervened state (\$Z=1\$), within the context where \$Y\_\{upt\}=0\$.} \\
\texttt{\$\$NIE = P(M=1 | Y=0, Z=1) {-} P(M=1 | Y=0, Z=0)\$\$} \\
\texttt{\$\$NIE = 51\textbackslash{}\% {-} 44\textbackslash{}\% = +7\textbackslash{}\%\$\$} \\
\\ 
\texttt{**Step 3: Conclusion**} \\
\texttt{The Natural Indirect Effect is \$+7\textbackslash{}\%\$. This positive value indicates that when \$Y\_\{upt\}\$ causes \$Z\_\{upj\}\$ to increase (from 0 to 1), the probability of \$M\_\{uvq\}\$ increases by 7 percentage points. Therefore, the effect of \$Y\_\{upt\}\$ on \$M\_\{uvq\}\$ through \$Z\_\{upj\}\$ is **positive**, not negative.} \\
\\ 
\texttt{The answer is **No**.} \\
            }
        }
    }
    & & \\ \\

    \theutterance \stepcounter{utterance}  
    & & & \multicolumn{2}{p{0.3\linewidth}}{
        \cellcolor[rgb]{0.9,0.9,0.9}{
            \makecell[{{p{\linewidth}}}]{
                \texttt{\tiny{[GM$|$GM]}}
                \texttt{No} \\
            }
        }
    }
    & & \\ \\

    \theutterance \stepcounter{utterance}  
    & & & \multicolumn{2}{p{0.3\linewidth}}{
        \cellcolor[rgb]{0.9,0.9,0.9}{
            \makecell[{{p{\linewidth}}}]{
                \texttt{\tiny{[GM$|$GM]}}
                \texttt{yes} \\
            }
        }
    }
    & & \\ \\

    \theutterance \stepcounter{utterance}  
    & & & \multicolumn{2}{p{0.3\linewidth}}{
        \cellcolor[rgb]{0.9,0.9,0.9}{
            \makecell[{{p{\linewidth}}}]{
                \texttt{\tiny{[GM$|$GM]}}
                \texttt{game\_result = LOSE} \\
            }
        }
    }
    & & \\ \\

\end{supertabular}
}

\end{document}
